% ===== PRÉAMBULE CANONIQUE — à coller VERBATIM en tête de CHAQUE .tex =====
\documentclass[11pt,a4paper]{article}
\usepackage[utf8]{inputenc}\usepackage[T1]{fontenc}\usepackage{lmodern}
\usepackage[french]{babel}
\usepackage{amsmath,amssymb,mathtools}\usepackage{icomma}
\usepackage[version=4]{mhchem}          % \ce{...} : équations & formules
\usepackage{siunitx}\sisetup{locale=FR} % \qty{}{}, \si{}, \ang{}
\usepackage{chemfig}                    % \chemfig{...} : structures développées
\usepackage{xcolor}\usepackage{colortbl}
\usepackage{tikz}\usepackage{pgfplots}\pgfplotsset{compat=1.18}
\usepackage[most]{tcolorbox}\usepackage{enumitem}\usepackage{array,booktabs}
\usepackage[a4paper,margin=2.2cm]{geometry}
\usetikzlibrary{arrows.meta,calc,angles,quotes,positioning,patterns,backgrounds,shapes.geometric}
\definecolor{primary}{RGB}{222,49,99}\definecolor{primarydark}{RGB}{150,22,55}
\definecolor{defcol}{RGB}{0,114,178}\definecolor{methcol}{RGB}{52,92,148}
\definecolor{excol}{RGB}{0,135,90}\definecolor{attncol}{RGB}{200,40,55}
\tcbset{boxbase/.style={enhanced,breakable,boxrule=0.9pt,arc=2.5pt,
  left=3mm,right=3mm,top=2.3mm,bottom=2.3mm,fonttitle=\bfseries,coltitle=white}}
% --- boîtes TUTORIEL (noms EXACTS, ne pas renommer) ---
\newtcolorbox{cle}[1][]{boxbase,colback=defcol!6,colframe=defcol,
  title={\textsc{À retenir}\ifx\relax#1\relax\relax\else\textmd{ : \itshape #1}\fi}}
\newtcolorbox{methode}[1][]{boxbase,colback=methcol!7,colframe=methcol,sharp corners,
  title={\textsc{Méthode}\ifx\relax#1\relax\relax\else\textmd{ --- \itshape #1}\fi}}
\newtcolorbox{exemplebox}[1][]{boxbase,colback=excol!5,colframe=excol,
  title={\textsc{Exemple}\ifx\relax#1\relax\relax\else\textmd{ : \itshape #1}\fi}}
\newtcolorbox{piege}[1][]{boxbase,colback=attncol!6,colframe=attncol,
  title={\textsc{Piège}\ifx\relax#1\relax\relax\else\textmd{ : \itshape #1}\fi}}
\newtcolorbox{remarque}[1][]{boxbase,colback=black!4,colframe=black!45,
  title={\textsc{Remarque}\ifx\relax#1\relax\relax\else\textmd{ : \itshape #1}\fi}}
% --- boîtes EXERCICE / CORRIGÉ (noms EXACTS, auto-comptées) ---
\newtcolorbox[auto counter]{exo}[1][]{boxbase,colback=primary!4,colframe=primary,
  title={\textsc{Exercice}~\thetcbcounter\ifx\relax#1\relax\relax\else\textmd{ --- \itshape #1}\fi}}
\newtcolorbox[auto counter]{corrige}[1][]{boxbase,colback=excol!4,colframe=excol,
  title={\textsc{Corrigé}~\thetcbcounter\ifx\relax#1\relax\relax\else\textmd{ --- \itshape #1}\fi}}
\newcommand{\R}{\mathbb{R}}\newcommand{\N}{\mathbb{N}}\newcommand{\Z}{\mathbb{Z}}\newcommand{\ds}{\displaystyle}
% (un \usepackage{fancyhdr}+en-tête est possible ; il est ignoré côté web)
% =========================================================================
\begin{document}
% THEME_TITLE: Identifier le réactif limitant et calculer l'excès
\sisetup{per-mode=symbol}

\begin{center}
{\Large\bfseries\color{primarydark} Thème 4 — Réactif limitant et réactif en excès}
\end{center}

\medskip
En pratique, on mélange rarement les réactifs dans les proportions exactes de l'équation. L'un d'eux
s'épuise le premier et \emph{stoppe} la réaction : c'est le \textbf{réactif limitant}. Tout le bilan
final se calcule à partir de lui.

\begin{cle}[limitant et excès]
Le \textbf{réactif limitant} est entièrement consommé le premier : sa disparition fixe les quantités
maximales de produits. Les autres réactifs, seulement en partie consommés, sont \textbf{en excès}
(il en \emph{reste} à la fin).
\end{cle}

\begin{methode}[le critère du rapport $n/\nu$]
Pour chaque réactif, on calcule $\dfrac{n_{\text{initial}}}{\nu}$ (quantité engagée divisée par son
coefficient). Le \textbf{plus petit rapport} désigne le réactif limitant :
\[
\boxed{\;\text{limitant}=\text{réactif de plus petit }\dfrac{n_i}{\nu_i}\;}
\qquad
\xi_{\max}=\min_i\!\left(\dfrac{n_i}{\nu_i}\right).
\]
Ensuite, \emph{tout} se calcule à partir du limitant : produits formés et excès restant.
\end{methode}

\begin{piege}[le limitant n'est pas « celui qu'il y a en plus petite quantité »]
On ne compare \textbf{jamais} les quantités brutes ! Pour $\ce{N2 + 3 H2 -> 2 NH3}$ avec
\qty{2}{\mole} de \ce{N2} et \qty{3}{\mole} de \ce{H2}, il y a \emph{plus} de \ce{H2}, et pourtant
$\frac{3}{3}=1<\frac{2}{1}=2$ : c'est \ce{H2} le limitant. C'est le \textbf{rapport} $n/\nu$ qui tranche.
\end{piege}

\begin{center}
\begin{tikzpicture}[font=\small]
  \def\u{34} % 1 unité de rapport -> 34 mm
  \draw[->,black!60] (0,0) -- (0,3.1) node[left,black,font=\scriptsize]{$\dfrac{n}{\nu}$ (\si{\mole})};
  % A : 0,05  -> limitant
  \fill[attncol!75] (0.8,0) rectangle (1.8,{0.05*\u});
  \node[below,font=\scriptsize] at (1.3,-0.05) {\ce{A} : $\frac{0,10}{2}$};
  \node[white,font=\scriptsize] at (1.3,{0.025*\u}) {$0,05$};
  \node[attncol,font=\scriptsize\bfseries,align=center] at (1.3,{0.05*\u+0.3}) {limitant};
  % B : 0,067 -> excès
  \fill[excol!70] (3.0,0) rectangle (4.0,{0.0667*\u});
  \node[below,font=\scriptsize] at (3.5,-0.05) {\ce{B} : $\frac{0,20}{3}$};
  \node[white,font=\scriptsize] at (3.5,{0.033*\u}) {$0,067$};
  \node[excol!60!black,font=\scriptsize\bfseries] at (3.5,{0.0667*\u+0.3}) {excès};
  \node[align=left,font=\scriptsize,black!70] at (7.4,1.5)
     {Le plus petit rapport ($0,05$)\\ est celui de \ce{A} : \ce{A} est\\ \textbf{limitant},
      \ce{B} est en excès.\\ Ici $\xi_{\max}=0,05$ \si{\mole}.};
\end{tikzpicture}
\end{center}

\begin{remarque}[et l'excès restant ?]
La quantité d'un réactif en excès qui \emph{reste} vaut : $n_{\text{restant}}=n_{\text{initial}}-
n_{\text{consommé}}$, où $n_{\text{consommé}}$ se calcule à partir du limitant par la règle de trois.
\end{remarque}

\end{document}
